\Needspace{6\baselineskip}
\Needspace{6\baselineskip}
\section{Publication terminale du continu unique par \(5\) lecteurs non permutables}
\Needspace{6\baselineskip}
\subsection{Factorisation de l’état enrichi commun en \(5\) publications terminales}
\Needspace{6\baselineskip}
\subsubsection{Dépendance causale des réalisations terminales à l’égard du continu généalogique intégralement construit}
Les cinq problèmes ne font pas appel à cinq structures mathématiques différentes. Ils n’obligent pas non plus à construire cinq chemins indépendants qui, après avoir progressé séparément, devraient être réunis artificiellement dans une conclusion commune. Cette image appartient à une étape antérieure de l’exposé. Le point de départ actuel est plus fort : avant que n’apparaissent les noms de Riemann, Navier–Stokes, Yang–Mills, Hodge ou Birch–Swinnerton–Dyer, il existe déjà un unique objet mathématique enrichi dont les opérations internes produisent simultanément les caractéristiques nécessaires aux cinq fermetures.\par
C’est pourquoi ce récit ne commence pas par les problèmes. Il commence par la construction du continu.\par
L’Holographie modulaire triadique ne reçoit pas le continu comme un arrière-plan préalablement donné. Elle ne part pas davantage de la droite réelle, d’un espace fonctionnel classique, d’une variété algébrique, d’une connexion de jauge ou d’une courbe elliptique. Elle part d’une arithmétique finie dotée d’orientation, de mémoire et d’une capacité de prolongement. APP fournit le substrat positionnel et arithmétique ; TRIT transforme chaque opération en une transition orientée ; TPK conserve la généalogie complète de ces transitions et permet de les prolonger sans effacer le chemin qui les a produites.\par
L’étape décisive tient à ce que la sortie d’une opération cesse d’être seulement une valeur. Chaque émission porte son feuillet arithmétique, son résidu, son quotient, son orientation, sa retenue, sa mémoire, sa frontière, son parcours, sa multisection, sa survie, son terme terminal et son lecteur. Le résultat n’est plus un point isolé, mais un état qui conserve sa provenance, ses possibilités de continuation et les relations à préserver lorsqu’il change d’échelle.\par
APP, TRIT et TPK ne forment pas trois explications superposées. Ils forment une seule séquence causale :\par
\[
\mathrm{APP}
\longrightarrow
\mathrm{TRIT}
\longrightarrow
\mathrm{TPK}
\longrightarrow
\text{état enrichi unique}.
\]
APP fixe le support arithmétique commun. TRIT introduit l’orientation effective de chaque transition, en distinguant avancée, seuil et ouverture future. TPK réunit l’information visible et l’information transportée : non seulement ce qu’une opération fournit, mais aussi la retenue qu’elle déplace, la frontière qu’elle modifie, la mémoire qu’elle conserve et les extensions qu’elle permet.\par
La connexion nonadique ne ramène pas le système à un état vide. Après un tour complet, la phase réapparaît, mais l’histoire n’est pas effacée. Le retour \(9\to1\) est un retour de phase, non une réinitialisation de l’état. Chaque tour ajoute profondeur, mémoire et compatibilité. Ainsi, la continuité ne s’obtient ni en répétant neuf filtres ni en traversant neuf canaux indépendants. Elle s’obtient en prolongeant une seule connexion complète dont les neuf segments sont des cartes internes d’une même holonomie.\par
\Needspace{6\baselineskip}
\subsubsection{Émergence corrélative de \texorpdfstring{\(5\)}{5} actions intrinsèques non exhaustives}
Une opération commune agit sur chaque état enrichi. Il n’existe pas d’abord cinq domaines, ni cinq sélections, ni cinq copies de l’état. La même émission est lue simultanément selon cinq caractéristiques internes désignées par les discriminants \(\mathrm{sp}\), \(\mathrm{sol}\), \(\mathrm{gau}\), \(\mathrm{coh}\) et \(\mathrm{det}\).\par
Ces symboles ne désignent pas des branches et n’épuisent pas les propriétés du continu. Ils désignent \(5\) actions intrinsèques du même objet dont la corrélation est décisive pour les réalisations terminales de cette Partie.\par
La caractéristique spectrale–polaire, \(sp\), enregistre l’expression orientée de l’état, son complément, son défaut, sa mémoire terminale et la relation entre les deux lectures spéculaires d’une même forme.\par
La caractéristique solénoïdale, \(sol\), enregistre la manière dont une transition conserve l’état tout en répartissant dissipation, advection, retenue, mémoire, microstructure et coquille sans les compter plusieurs fois.\par
La caractéristique de jauge, \(gau\), enregistre l’équivalence locale, l’incidence, l’holonomie, la courbure et le passage compatible au quotient physique.\par
La caractéristique cohomologique, \(coh\), enregistre les extensions entre niveaux, les obstructions relatives, la compatibilité des relèvements et la survie d’une classe à travers tout raffinement.\par
La caractéristique déterminantale, \(det\), enregistre l’incidence finie-globale, les pages de filtration, les complexes locaux, leurs droites déterminantes et le transport d’une base arithmétique jusqu’à une expression analytique.\par
Mais aucune de ces propriétés n’existe seule. Toutes agissent sur les mêmes émissions et conservent la même épaisseur généalogique. La fibre qui transporte la caractéristique spectrale est la même fibre enrichie dont l’orientation participe au bilan solénoïdal, dont la frontière porte la lecture de jauge, dont le prolongement soutient l’extension cohomologique et dont l’incidence entre dans la droite déterminante.\par
On n’engendre pas cinq objets pour les regrouper ensuite. On engendre un unique objet qui possède indissociablement ces cinq capacités.\par
Une vue ultérieure n’est donc pas une projection qui crée une branche. Elle est seulement une manière d’exposer, pour une application concrète, une caractéristique qui agissait déjà au sein de la totalité. La vue ne sélectionne pas rétrospectivement des états favorables et ne conserve pas l’entrée comme première coordonnée pour la retrouver ensuite de manière tautologique. Elle rend visible une structure produite et conservée tout au long de la généalogie.\par
\Needspace{6\baselineskip}
\subsubsection{Antériorité causale du continu déjà construit comme limite du prolongement compatible}
À chaque profondeur existe un état fini, mais complet à son niveau. Les extensions entre niveaux ne peuvent modifier arbitrairement l’information antérieure. Elles doivent respecter les fibres, les relations d’extension, les nombres, l’orientation, la retenue, la mémoire, la frontière, le parcours, la multisection, la survie, le terme terminal et les lecteurs.\par
La naturalité exige qu’opérer puis tronquer donne le même résultat que tronquer puis opérer. Grâce à cette compatibilité, les histoires finies ne demeurent pas des approximations déconnectées. Elles forment un système inverse. Le continu apparaît lorsqu’on considère la famille complète des histoires compatibles, et non lorsqu’on postule d’emblée une collection infinie de points.\par
Le continu HMT est donc mémoire cohérente de chemins. Un point sans généalogie serait une ablation : il conserverait la valeur visible et supprimerait précisément ce qui permet de démontrer comment elle survit à tout raffinement.\par
Ce n’est qu’après avoir établi cette compatibilité qu’apparaît la structure discrète complète du continu. Ce n’est qu’ensuite qu’une réalisation archimédienne peut être exprimée. La droite réelle n’est pas la matière première de la construction ; elle est une sortie de l’un de ses lecteurs terminaux. Il en va de même pour les domaines classiques des cinq problèmes : ils apparaissent lorsque le continu est déjà construit et qu’une interface ultérieure traduit l’une de ses caractéristiques intrinsèques dans le langage de l’objectif correspondant.\par
La séquence causale correcte est donc :\par
\[
\begin{aligned}
\mathrm{APP}
&\longrightarrow \mathrm{TRIT}
\longrightarrow \mathrm{TPK}
\longrightarrow \text{état enrichi unique}\\
&\longrightarrow 04\mathrm b1
\longrightarrow 04\mathrm b2
\longrightarrow 04\mathrm b3\\
&\longrightarrow \text{actions intrinsèques corrélatives}
\longrightarrow \text{terminal et limite}
\longrightarrow \mathfrak C_{\mathrm{cont}}^{\mathrm{disc}}.
\end{aligned}
\]
Les entrelaceurs spécifiques interviennent ensuite, et pas avant :\par
\[
\mathfrak C_{\mathrm{cont}}^{\mathrm{disc}}
\longrightarrow
\begin{cases}
\text{forme explicite de Weil},\\
\text{évolution tridimensionnelle incompressible},\\
\text{théorie de jauge quantique},\\
\text{publication des classes de Hodge},\\
\text{comparaison déterminantielle BSD}.
\end{cases}
\]
Ces cinq traductions terminales ne définissent pas la structure qu’elles utilisent. Le problème classique apporte ses données après coup : une fonction test pour Riemann, des données initiales et une force pour Navier–Stokes, un groupe compact simple pour Yang–Mills, une variété et une classe pour Hodge, ou une courbe elliptique pour BSD. Aucune de ces données ne participe à la génération préalable du continu.\par
La direction causale est ainsi protégée. L’objectif ne sélectionne pas l’état qui devra le résoudre. L’état existe avant l’objectif ; l’entrelaceur démontre seulement que la structure interne et la formulation classique expriment la même opération dans deux langages.\par
\Needspace{6\baselineskip}
\subsection{1re résolution : Riemann et la positivité comme complément construit}
L’hypothèse de Riemann est généralement présentée comme une assertion concernant la position des zéros non triviaux de la fonction zêta. Dans cette formulation, le problème semble demander d’examiner une population de zéros déjà donnée et de démontrer qu’ils s’alignent tous sur la droite critique.\par
La lecture HMT inverse cette perspective. Elle ne commence pas par demander où sont les zéros. Elle demande quelle structure engendre simultanément les deux membres de la formule explicite et ce qui empêche leur différence de devenir négative.\par
La caractéristique spectrale–polaire du continu fournit une double expression du même état. Une orientation recueille le canal positif ; l’autre recueille son image spéculaire. Toutes deux contiennent, sous un régulateur commun, la contribution première, la contribution gamma–scalaire et la contribution polaire. Il ne s’agit ni de trois preuves différentes ni de trois sommes accolées à la fin. Ce sont trois lignes d’une même colligation construite à partir des données initiales.\par
Le point essentiel est que cette colligation est construite avant d’invoquer la positivité qu’elle devra démontrer. Son entrée, ses sorties, sa perte et son état terminal sont typés dans le même espace enrichi. L’identité qui relie les deux expressions n’est pas introduite comme une égalité souhaitée entre formules classiques ; elle naît de la dynamique interne de la connexion.\par
Entre l’expression d’entrée et l’expression spéculaire peut apparaître, à profondeur finie, un résidu de sortie. La structure nonadique lui impose une récurrence antispéculaire. À chaque tour complet, le résidu suivant est transporté isométriquement, mais multiplié par le facteur\par
\[
q=\frac{5}{9}.
\]
Si \(\Omega_n\) représente ce défaut à la profondeur \(n\), la récurrence prend la forme\par
\[
\Omega_n=q\,V_n\Omega_{n+1},
\]
avec \(V_n\) isométrique et la famille coinductivement bornée. En itérant \(N\) fois,\par
\[
\|\Omega_n\|
\le q^N\sup_{m\ge n}\|\Omega_m\|.
\]
Comme \(0<q<1\), le membre de droite tend vers zéro. Le résidu ne s’annule donc pas parce qu’il a été omis ou parce qu’une identification entre les deux expressions a été supposée. Il s’annule parce que la récurrence complète le contracte à toutes les profondeurs compatibles. La sortie spéculaire coïncide alors avec l’expression négative construite.\par
Une fois ce lien établi, l’isométrie construite à partir des données initiales fournit l’identité d’énergie. À profondeur finie, la norme de l’expression positive se décompose en trois parties : l’expression négative, la perte accumulée et l’état terminal. Le terme terminal est explicitement conservé. Il doit être conservé jusqu’au passage à la limite, car il contient la mémoire de la troncature et certifie que le télescopage prolonge l’identité finie par une convergence démontrée.\par
La décomposition finie de la norme satisfait l’identité exacte :\par
\[
\|\text{publication positive}\|^2
=
\|\text{publication négative}\|^2
+
\|\text{pertes jusqu'à }N\|^2
+
q^{2N}\|\text{terminal}\|^2.
\]
Chaque terme possède une origine concrète. L’expression négative est la sortie spéculaire déjà identifiée. Les pertes réunissent exactement le complément qui n’a pas été exprimé par cette orientation. Le terme terminal conserve ce qui demeure encore au-delà de la coupe.\par
Ce n’est qu’après avoir écrit l’identité complète que l’on fait \(N\to\infty\). Alors, et alors seulement, le terme terminal disparaît parce que \(q^{2N}\to0\). La différence entre les deux expressions est transformée en carré d’un opérateur de perte :\par
\[
\mathcal I_{\mathrm{GW}}
=
E_R^{*}E_R.
\]
C’est la forme de Guinand–Weil. Sa positivité n’est plus une conjecture ajoutée à la construction :\par
\[
\langle f,\mathcal I_{\mathrm{GW}}f\rangle
=
\|E_Rf\|^2
\ge0.
\]
La structure HMT n’obtient pas cette inégalité en dissimulant la partie négative, en restreignant le domaine aux fonctions qui la satisfont déjà, ou en introduisant un projecteur dont l’existence serait équivalente à la conclusion. Elle construit le complément et démontre que la forme est exactement son carré.\par
L’expression commune est ensuite matérialisée. Un même encodeur conserve les canaux premier, gamma–scalaire et polaire. Un projecteur interne produit l’un des deux moments ; son complément produit l’autre. La différence des moments n’est pas une comparaison externe entre formules semblables : c’est la norme de la composante orthogonale construite par la colligation.\par
Lorsqu’on élargit la fenêtre, on n’oblige pas les colonnes individuelles à satisfaire de fausses isométries entre troncatures distinctes. Ce qui est conservé de manière cofinale, c’est la forme, le complément atteignable et la naturalité de son expression. Les incréments introduits par l’élargissement de la coupe apparaissent dans les deux orientations avec la même contribution neutre et ne modifient pas la différence directrice. On obtient ainsi une famille compatible de formes positives sur une classe séparante.\par
La dernière transition est le critère réversible de Weil. La théorie classique affirme que la positivité de la forme explicite, sur la classe adéquate de fonctions tests, équivaut au fait que tous les zéros non triviaux se trouvent sur la droite critique. HMT fournit précisément cette positivité, mais comme conséquence d’une identité opératorielle construite à l’intérieur du continu.\par
\Needspace{5\baselineskip}
Por tanto,\par
\[
\operatorname{Re}\rho=\frac12
\]
pour tout zéro non trivial \(\rho\) de la fonction zêta.\par
La droite critique apparaît ici de deux manières concordantes. Géométriquement, elle est le lieu fixe de l’involution qui échange les deux orientations spectrales. Du point de vue opératoriel, elle est l’unique position compatible avec la positivité du complément de l’expression. La symétrie explique pourquoi la droite centrale est distinguée ; l’identité quadratique démontre pourquoi les zéros ne peuvent quitter cette droite.\par
La résolution ne consiste pas à observer que l’équation fonctionnelle possède une symétrie. Elle consiste à construire l’objet dont le défaut mesure la rupture de cette symétrie, à démontrer que le défaut s’annule coinductivement et à transformer la différence restante en une norme. C’est la combinaison de la généalogie, du terme terminal, de la contraction, de l’expression commune et de la positivité qui clôt le problème.\par
La première des cinq résolutions terminales est ainsi établie. Non comme une branche indépendante, mais comme la traduction classique de la caractéristique spectrale–polaire que le continu unique possédait dès sa construction.\par

\Needspace{6\baselineskip}
\section{Fermetures dynamiques du continu : transport solénoïdal et quotient de jauge}
\Needspace{6\baselineskip}
\subsection{Dynamique et courbure du continu généalogique}
La première résolution a montré comment la caractéristique spectrale–polaire du continu s’exprime comme positivité de Weil. Les deux résolutions suivantes interrogent le même état sous deux autres aspects consubstantiels : sa capacité à transporter une dynamique non linéaire sans perdre énergie, mémoire ni régularité, et sa capacité à soutenir une théorie de jauge locale dont le vide est séparé par un saut spectral positif.\par
Navier–Stokes et Yang–Mills ne reçoivent pas deux structures nouvelles. Ils reçoivent deux entrelaceurs différents à partir du même continu déjà construit. Dans le premier cas, l’entrelaceur exprime l’état comme un champ de vitesse incompressible avec viscosité, advection et pression. Dans le second, il l’exprime comme un réseau de jauge raffinable avec holonomie, courbure, mesure euclidienne et Hamiltonien physique.\par
La différence entre les deux problèmes ne rompt pas l’unité de la construction. Pour Navier–Stokes, le point crucial est de démontrer qu’aucun transfert non linéaire ne disparaît hors du bilan et que la régularité obtenue est uniforme lorsqu’on retire les coupures. Pour Yang–Mills, le point crucial est de démontrer que le saut fini appartient à la même théorie de jauge qui atteint le continu, qu’il subsiste au quotient physique et qu’il agit dans le secteur de courbure, non dans un facteur auxiliaire.\par
\Needspace{6\baselineskip}
\subsection{Existence, régularité et unicité globales pour Navier–Stokes par borne uniforme et compacité}
Le problème tridimensionnel de Navier–Stokes ne naît pas de la seule présence d’un terme non linéaire. La difficulté apparaît parce que l’expression classique du champ de vitesse tend à masquer le lieu où se stocke le transfert entre échelles. La viscosité est visible. L’énergie cinétique est visible. L’advection apparaît également dans l’équation. Mais la comptabilité complète de l’interaction entre profondeur, mémoire, retenue, microstructure et coquille est comprimée en une seule fonction visible.\par
C’est précisément cette compression que la caractéristique solénoïdale du continu évite.\par
L’état enrichi ne conserve pas uniquement le champ \(u\). Il conserve simultanément le champ de Galerkin, l’orientation temporelle, les mémoires positive et signée, la retenue entre coupures, la composante microscopique, la coquille des fréquences nouvelles, le parcours, l’archive de l’évolution et le terme terminal. L’équation classique est une expression de cet état ; elle n’en est pas la définition complète.\par
Les données initiales \(u_0\), la force \(f\), la viscosité \(\nu>0\) et la régularité \(s>5/2\) étant fixées, la projection de Galerkin produit\par
\[
\partial_tu_M+\nu A_Mu_M+B_M(u_M,u_M)=f_M.
\]
Mais \(M\) ne sélectionne pas une solution favorable. Elle fixe seulement une résolution finie dans laquelle toutes les opérations peuvent être écrites exactement. L’objectif est de construire une estimation indépendante de \(M\) et de la profondeur nonadique \(J\). Seule une borne présentant cette uniformité peut passer au continu.\par
\Needspace{6\baselineskip}
\subsubsection{Décomposition de la mémoire visible et de la mémoire généalogique dans l’équation d’évolution}
L’interaction non linéaire possède un signe. À certains instants, elle apporte de l’énergie à une échelle ; à d’autres, elle lui en retire. Si l’on ne conserve que sa valeur signée, l’amplitude totale transférée peut être perdue. Si l’on ne conserve que sa valeur absolue, l’orientation est perdue.\par
La structure maintient les deux.\par
Une mémoire \(D_M\) exprime le signe du transfert advectif. Une autre mémoire \(H_M\) conserve son amplitude positive. Leur séparation permet d’enregistrer simultanément\par
\[
\dot D_M=\mathscr B_{s,M},
\qquad
\dot H_M=\frac{1+r}{r-1}|\mathscr B_{s,M}|,
\]
où \(r=\pi_{\rm HMT}/\varphi_{\rm HMT}^2>1\).\par
Cela n’ajoute pas d’énergie fictive. Cela dédouble deux aspects distincts d’un seul transfert : orientation et coût. Les confondre détruirait de l’information ; les additionner sans orthogonalité doublerait l’énergie. La construction montrera explicitement qu’aucune de ces deux déformations ne se produit.\par
\Needspace{6\baselineskip}
\subsubsection{Indépendance de la donnée initiale à l’égard de l’histoire généalogique ultérieure}
L’objection la plus sérieuse à toute fermeture opératorielle de Navier–Stokes serait que la borne uniforme ait été introduite rétrospectivement : on choisirait d’abord une solution régulière, puis on coderait son histoire, pour annoncer enfin que le code contrôle cette même solution.\par
La construction empêche cette circularité en fixant l’espace source avant de relever l’évolution :\par
\[
\mathcal D_T
=
H^s_\sigma(\mathbb T^3)
\times
L^2(0,T;H^{s-1}_\sigma(\mathbb T^3)),
\qquad
d=(u_0,f).
\]
La racine \(J_T^{\rm sol}d\) dépend exclusivement de \(u_0\) et de \(f\). Elle ne reçoit en entrée ni \(u_M\), ni la perte future, ni le terme terminal, ni la norme à borner.\par
La force n’est pas non plus insérée comme un port additif artificiel. L’équation de Carleman montre qu’elle agit de manière mixte : chaque occurrence de \(f\) se combine avec un degré antérieur de l’état. Le relèvement utilise donc une création mixte état–force. La signature chronologique de \(f\) enregistre tous ses temps ordonnés, tandis que la tour PC–Fock enregistre tous les degrés non linéaires de l’état.\par
La composition reproduit matériellement la série de Dyson du générateur non autonome. L’application du lecteur normalisé de degré un restitue exactement \(u_M(t)\). Elle ne restitue ni une approximation symbolique ni une variable ressemblant à la solution : elle restitue la solution de l’EDO de Galerkin, car toutes deux satisfont la même équation finie et la même donnée initiale.\par
La construction causale à partir des données initiales suit donc l’ordre\par
\[
(u_0,f)
\longrightarrow
J_T^{\rm sol}(u_0,f)
\longrightarrow
\text{tour état–force}
\longrightarrow
\text{histoire de Galerkin}.
\]
L’histoire généalogique est construite comme sortie ultérieure des données initiales et est exclue de leur définition initiale.\par
\Needspace{6\baselineskip}
\subsubsection{L’identité de diffusion qui conserve la puissance croisée}
La force et la dissipation sont organisées au moyen de deux ondes :\par
\[
a_T(f)
=
\frac{1}{2\sqrt\nu}A^{-1/2}\Lambda^sP_\sigma f,
\qquad
z_M(u)=\sqrt\nu A_M^{1/2}\Lambda^su,
\]
et de l’onde résiduelle\par
\[
b_M=z_M-P_Ma_T(f).
\]
La puissance appliquée n’est pas remplacée par la norme de la force. Elle est conservée exactement comme terme croisé :\par
\[
\mathscr F_{s,M}
=
2\operatorname{Re}\langle z_M,P_Ma_T(f)\rangle.
\]
La complétion du carré donne le bilan exact\par
\[
\dot G_{M,j,T}
+\|b_M\|^2
+\lvert\mathscr B_{s,M}\rvert
=
\|P_Ma_T(f)\|^2.
\]
C’est l’un des points qui rendent la fermeture opérante. Si l’on remplaçait \(b=z-a\) par \(z\), le terme croisé représentant le travail de la force disparaîtrait et l’identité serait fausse. La structure ne maîtrise pas la non-linéarité en l’écartant : elle la transforme en un canal positif orienté et conserve en outre l’interférence exacte entre force et dissipation.\par
\Needspace{6\baselineskip}
\subsubsection{Factorisation de \(6\) pertes par un opérateur unique sans multiplicité parasite}
L’histoire complète distingue six sorties :\par
\begin{itemize}
\item dissipation ;
\item advection orientée ;
\item retenue entre coupures ;
\item mémoire microscopique ;
\item coquille de fréquences nouvelles ;
\item mémoire non héritée.
\end{itemize}
Ce ne sont pas six énergies ajoutées à un bilan préalable. Ce sont six lignes orthogonales d’un unique opérateur terminal de perte. L’identité de Parseval établit\par
\[
\|L_{M,j,T}^{\rm term}x\|^2
=
\sum_{\mathfrak a}
\|L_{M,j,T}^{\mathfrak a}x\|^2.
\]
L’orthogonalité n’est pas simplement nominale. Les profondeurs occupent des intervalles temporels disjoints ; retenue et coquille vivent dans des plages spectrales distinctes ; la microstructure appartient au complément \(Q_{\rm micro}\) ; la mémoire nouvelle est séparée de la mémoire héritée ; et les deux feuillets de l’advection ne peuvent être simultanément positifs en un même point.\par
La comptabilité ne double donc pas l’interaction. La somme des six normes est exactement la norme de l’opérateur commun. Si l’une d’elles n’était pas une projection de cet opérateur, Parseval échouerait avant l’obtention de la borne. Le formalisme contient ainsi son propre réfutateur.\par
\Needspace{6\baselineskip}
\subsubsection{Persistance de l’objet terminal sous le passage à la limite}
Toute troncature conserve quelque chose qui n’a pas encore été exprimé. Ce reste est le terme terminal. Il contient l’énergie visible à la fin de l’intervalle et la mémoire orientée qui demeure entre le temps observé et l’horizon \(T\).\par
La colonne terminale est construite en premier. Son Gram est formé ensuite. On ne définit pas une métrique souhaitée pour inventer ultérieurement un opérateur qui la réalise. La colonne contient matériellement la sortie terminale et toutes les pertes transportées :\par
\[
\mathbf G_{M,j;J,T}^{\rm term}
=
\begin{bmatrix}
\text{terminal transporté}\\
\text{perte}_{J-1}\\
\vdots\\
\text{perte}_{j}
\end{bmatrix}.
\]
Sa récurrence par blocs produit l’identité de Stein\par
\[
U_{M,j,T}
=
\Gamma_{M,j,T}^{*}
U_{M,j+1,T}
\Gamma_{M,j,T}
+
(L_{M,j,T}^{\rm term})^{*}
L_{M,j,T}^{\rm term}.
\]
Cette identité ne définit pas rétroactivement la norme de l’état. Elle s’obtient en multipliant une colonne déjà construite par son adjointe. Pour la même raison, le télescopage global est une identité entre opérateurs existants, et non une hypothèse de stabilité.\par
La mémoire future n’est pas non plus introduite dans la racine initiale. Elle réside dans la sortie terminale. Lorsqu’on prend l’estimation intégrale en \(\tau=T\), son support temporel se vide exactement. Elle n’est pas effacée par commodité.\par
\Needspace{6\baselineskip}
\subsubsection{Borne uniforme préalable et compacité d’Aubin–Lions}
La racine, l’encodeur de niveau, la colonne et l’isométrie commune satisfont la factorisation dérivée des données initiales\par
\[
\mathcal G_{M,J,T}^{\rm sol}
\Lambda_{M,T}
J_T^{\rm sol}d
=
I_{M,J}J_T^{\rm sol}d,
\qquad
I_{M,J}^{*}I_{M,J}=I.
\]
Le membre de gauche est l’histoire complète à la coupure \((M,J)\). Celui de droite dépend de l’unique racine construite à partir des données. Puisque \(I_{M,J}\) est isométrique, la norme de l’histoire est dominée par la norme source avec une constante indépendante de \(M\) et de \(J\).\par
C’est l’étape qui résout la difficulté. L’uniformité ne provient pas d’une estimation de dimension finie dont la constante croît avec la coupure. Elle ne provient pas non plus d’une définition de l’espace des histoires au moyen de la perte même que l’on veut contrôler. La racine existe avant la coupure et toutes les coupures en sont des réalisations isométriques.\par
La colonne terminale traduit cette borne structurelle en grandeurs physiques :\par
\[
\begin{aligned}
&\sup_M\sup_{0\le t\le T}\|u_M(t)\|_{H^s}^2
+\nu\sup_M\int_0^T\|u_M(t)\|_{H^{s+1}}^2\,dt\\
&\quad
+\sup_M\int_0^T|\mathscr B_{s,M}(u_M(t))|\,dt
+\text{carry}+\text{micro}+\text{coquille}+\text{mémoire}
\le C_T(u_0,f).
\end{aligned}
\]
La constante est finie pour tout \(T<\infty\) et ne dépend ni de \(M\) ni de \(J\). L’inégalité ne contrôle pas seulement le champ visible. Elle contrôle simultanément le travail transporté et les secteurs qu’une expression classique tend à oublier.\par
\Needspace{6\baselineskip}
\subsubsection{Obstruction de signe dans la factorisation de Kalman–Yakubovich–Popov}
Faire dépendre la fermeture globale d’une factorisation KYP portant un signe incorrect introduit un carré positif et prétend le retrouver ensuite avec un signe négatif. Cette égalité est algébriquement fausse.\par
La démonstration exclut cette identité et construit directement l’estimation uniforme nécessaire.\par
La carte KYP conservée est une identité finie correcte sous ses propres conditions. Elle sert de contrôle auxiliaire de chaque coupure, mais ne fournit pas l’uniformité lorsqu’on retire \(M\) et \(J\). La preuve globale n’utilise pas de borne de \(Q^{-1}\), ne prend pas la limite à l’intérieur de KYP et ne dépend pas de ce signe.\par
L’obstruction du bloc KYP portant un signe incorrect n’affecte donc pas la démonstration. L’uniformité est établie par la construction directe à partir des données, avec sa colonne terminale et son identité de Stein.\par
\Needspace{6\baselineskip}
\subsubsection{Convergence de la borne uniforme vers la solution classique globale}
La borne précédente fournit\par
\[
u_M
\quad\text{borné dans}\quad
L^\infty(0,T;H^s)
\cap L^2(0,T;H^{s+1}).
\]
Comme \(s>5/2\), le produit de Sobolev contrôle la non-linéarité :\par
\[
B:H^s_\sigma\times H^s_\sigma
\longrightarrow H^{s-1}_\sigma.
\]
L’équation de Galerkin fournit en outre une borne uniforme de \(\partial_tu_M\) dans \(L^2(0,T;H^{s-1})\). Les inclusions\par
\[
H^{s+1}\Subset H^s\hookrightarrow H^{s-1}
\]
permettent d’appliquer Aubin–Lions. On obtient une convergence forte dans \(L^2H^s\), suffisante pour transporter le terme bilinéaire :\par
\[
B_M(u_M,u_M)\longrightarrow B(u,u)
\quad\text{dans}\quad
L^1(0,T;H^{s-1}).
\]
La limite satisfait l’équation incompressible. La pression n’est pas introduite comme une autre donnée cachée ; elle est reconstruite ensuite au moyen de\par
\[
-\Delta p=\partial_i\partial_j(u_i u_j),
\qquad
\int_{\mathbb T^3}p=0.
\]
L’unicité s’obtient sur la différence de deux solutions. Une estimation dans \(H^{s-1}\), suivie de Grönwall, impose que la différence soit nulle. Cela élimine la dépendance à l’égard de la sous-suite extraite par compacité.\par
La construction vaut pour tout horizon fini. Les solutions obtenues sur \([0,n]\) et \([0,n+1]\) coïncident sur leur chevauchement par unicité. Elles se recollent donc en une seule solution définie sur \([0,\infty)\).\par
La régularité lisse s’obtient sur toute l’échelle \(s+n\), à l’aide des mêmes lecteurs structurels ; l’unicité identifie toutes les limites et produit une solution globale lisse pour toute donnée périodique lisse. La borne requise est construite sur chaque horizon fini et pour chaque indice de Sobolev à partir de la donnée solénoïdale complète, comme le démontre le développement de référence \cref{thm:pdf2-ns-global}.\par
Enfin, la condition de moyenne nulle est une normalisation du secteur solénoïdal. Le changement de Galilée dépendant du temps sépare la moyenne et transporte bijectivement toute donnée périodique lisse vers le secteur de moyenne nulle, en préservant la divergence, la régularité lisse et les normes de Sobolev.\par
\Needspace{6\baselineskip}
\subsubsection{Existence, régularité et unicité globales de Navier--Stokes}
La caractéristique solénoïdale du continu, exprimée par son entrelaceur spécifique, produit une unique solution globale lisse des équations tridimensionnelles périodiques de Navier--Stokes pour toute donnée périodique lisse. La construction de Leray--Hopf est contenue comme étape d’énergie dans cette conclusion forte.\par
La défense de cette conclusion ne repose pas sur l’affirmation qu’il existe une structure capable de résoudre le problème. Elle repose sur cinq faits matériels :\par
la racine utilise uniquement \((u_0,f)\) ; la non-linéarité est reproduite par création mixte ; les six pertes sont orthogonales et ne sont pas additionnées deux fois ; le terme terminal demeure jusqu’au télescopage ; et la borne uniforme précède l’argument classique de compacité.\par
Aubin–Lions, la pression, Grönwall et le recollement ne sont pas la source de la régularité. Ils constituent la reconnaissance classique finale d’une régularité dont la structure a déjà empêché la perte.\par
\Needspace{6\baselineskip}
\subsection{3e résolution : Yang–Mills et le saut du secteur de courbure}
Le problème de Yang–Mills exige deux choses indissociables : construire une théorie quantique non triviale sur \(\mathbb R^4\) pour tout groupe compact, connexe et simple du domaine considéré, et démontrer que son Hamiltonien possède un saut de masse positif.\par
Il ne suffit pas de présenter un opérateur fini à spectre positif. Il ne suffit pas non plus d’écrire une action de Wilson, une mesure produit ou une observable ayant une valeur propre non nulle. Le saut doit appartenir à la même théorie de jauge qui possède localité, positivité par réflexion, covariance euclidienne, reconstruction de Wightman et courbure physique.\par
La caractéristique de jauge du continu conserve déjà le réseau raffinable, la mesure, les quatre réflexions, l’holonomie, la couleur, l’incidence, la retenue de fusion, le parcours, la mémoire, le terme terminal et le lecteur de courbure. Le groupe \(G\) et le couplage \(g\) interviennent ensuite, par l’entrelaceur spécifique. Ils ne sélectionnent pas rétrospectivement la généalogie.\par
\Needspace{6\baselineskip}
\subsubsection{Inéquivalence entre la fibre simple et la dynamique produit}
La première distinction dans l’argumentaire sépare deux générateurs qu’une ancienne notation pouvait confondre.\par
Le semi-groupe d’une seule fibre ne possède que les niveaux \(0\) et \(3\kappa_{\rm av}\). Le générateur produit complet intègre en revanche toutes les coordonnées indépendantes de l’état :\par
\[
D_m^{\rm prod}
=
3\kappa_{\rm av}
\sum_{\iota\in\mathcal I_m}(I-E_{m,\iota}),
\]
où les espérances \(E_{m,\iota}\) sont orthogonales et commutent.\par
Son spectre contient\par
\[
0,\quad
3\kappa_{\rm av},\quad
6\kappa_{\rm av},\quad
9\kappa_{\rm av},\ldots
\]
selon le nombre de coordonnées non constantes. Remplacer ce générateur par celui d’une fibre simple effacerait les multiplicités de Hoeffding et serait immédiatement falsifiable : un vecteur dépendant de deux facteurs doit avoir l’énergie \(6\kappa_{\rm av}\), et non \(3\kappa_{\rm av}\).\par
Le saut sera démontré pour le générateur produit, puis transporté vers le secteur physique. Il ne s’obtient pas en confondant deux dynamiques différentes.\par
\Needspace{6\baselineskip}
\subsubsection{Absence de coordonnées cachées et exclusion des modes nuls parasites}
L’état complet se factorise au moyen d’un indice exhaustif. Y apparaissent une seule fois :\par
\begin{itemize}
\item repères ;
\item liens ;
\item incidences ;
\item marques ;
\item ponts de raffinement ;
\item innovations nouvelles de chaque niveau.
\end{itemize}
Les sous-queues infinies sont enregistrées au moyen de leurs cylindres scalaires. Aucune coordonnée généalogique ne reste hors de la famille des espérances.\par
La décomposition de Hoeffding écrit l’espace comme somme orthogonale de secteurs \(\mathcal H_{m,S}\). Sur chacun d’eux,\par
\[
D_m^{\rm prod}u_S
=
3\kappa_{\rm av}|S|u_S.
\]
Le secteur \(S=\varnothing\) est l’intersection des images de toutes les espérances et se compose exactement des constantes :\par
\[
\bigcap_{\iota\in\mathcal I_m}
\operatorname{Ran}E_{m,\iota}
=
\mathbb C\Omega_m.
\]
Le vide est donc simple et tout vecteur orthogonal au vide a une énergie au moins égale à \(3\kappa_{\rm av}\).\par
Cela exclut l’objection selon laquelle le saut proviendrait de l’omission d’une coordonnée dans l’énumération. S’il manquait un facteur, il existerait des fonctions non constantes invisibles pour toutes les espérances et des modes nuls supplémentaires apparaîtraient. L’exhaustivité de l’indice empêche précisément cet échec.\par
\Needspace{6\baselineskip}
\subsubsection{Réalisation observable du saut spectral}
Une borne inférieure ne suffit pas à fixer la valeur exacte du saut. Un témoin atteignant le seuil est nécessaire.\par
Dans chaque collier d’incidence, la variable\par
\[
W_c(q)
=
\mathbf1_{\{q_1+\cdots+q_6=0\}}
-\frac13
\]
a une moyenne nulle et satisfait\par
\[
D_m^{\rm prod}W_c
=
3\kappa_{\rm av}W_c.
\]
Mais il reste à démontrer que \(W_c\) appartient au secteur de jauge physique.\par
L’action d’incidence s’écrit au moyen de la matrice\par
\[
(B^+x)_{ij}=x_i+x_j.
\]
Chaque composante de \(x\in\mathbb F_3^4\) apparaît trois fois dans la somme des six incidences. Dans \(\mathbb F_3\),\par
\[
\sum_{i<j}(B^+x)_{ij}
=
3\sum_i x_i
=
0.
\]
La condition définissant \(W_c\) est donc invariante de jauge. La moyenne sur le groupe de jauge complet conserve \(W_c\) ; elle ne l’exclut pas de l’espace physique. Ainsi, \(W_c\) est un vecteur propre physique et la valeur \(3\kappa_{\rm av}\) n’est pas seulement une borne : c’est le premier niveau spectral atteint.\par
\Needspace{6\baselineskip}
\subsubsection{Construction de la mesure préalable à l’estimation du saut}
Une arête grossière se raffine en neuf incréments ordonnés. La mesure fine n’est pas inventée en recopiant neuf fois la mesure grossière. Elle est construite au moyen du pont de chaleur conditionné par le produit des neuf incréments. Les temps des sous-segments peuvent être distincts et satisfont seulement à la condition que leur somme soit le temps grossier.\par
La convolution du noyau de chaleur garantit que le pont est une probabilité. Une carte triangulaire sépare alors la variable grossière des innovations nouvelles. La projection vers le niveau précédent multiplie les neuf incréments et oublie exclusivement les innovations.\par
De cette manière,\par
\[
(\widehat\pi_{m+1,m})_*
\widehat\mu_{m+1}
=
\widehat\mu_m.
\]
Les arêtes partagées par plusieurs faces demeurent une seule et même variable. Elles ne sont pas remplacées par des copies indépendantes. La retenue et la mémoire empêchent de recompter au niveau fin la face déjà présente au niveau grossier.\par
Sur cette famille projective, on construit d’abord le lecteur exact de \(F^2\), sa racine positive, l’action locale et la densité de Gibbs dépendant du couplage :\par
\[
d\widehat\mu_{m,g}^{\rm YM}
=
Z_{m,g}^{-1}
e^{-S_{m,g}}\,
d\widehat\mu_m^0.
\]
Les équations de Schwinger–Dyson et de Ward sont calculées pour cette mesure. Le saut n’est utilisé ni pour choisir \(g\), ni pour définir \(S_{m,g}\), ni pour repondérer la loi. Il apparaît ensuite, lors de l’étude du générateur réversible associé à la même mesure.\par
C’est la protection contre la circularité action–loi–saut : on construit d’abord l’action et sa mesure ; on démontre ensuite son spectre.\par
\Needspace{6\baselineskip}
\subsubsection{Descente du générateur au quotient de jauge}
L’espace complet comprend liens, repères, incidences, couleur, orientation, parcours, mémoire et terme terminal. La marginale qui ne conserve que les liens est une ablation ultérieure. Elle ne peut servir à nier une observable située dans la fibre d’incidence.\par
La projection physique s’obtient par moyenne de l’action du groupe de jauge complet. Puisque le générateur intègre les coordonnées de manière équivariante,\par
\[
[\mathbb P_m^{\rm phys},D_m^{\rm prod}]=0.
\]
Le semi-groupe descend donc au quotient physique. Son noyau est défini d’abord sur l’état complet, puis seulement sur les classes de jauge. On ne prend pas un opérateur restreint à une fibre pour le déclarer noyau sur tout l’espace.\par
Le vide simple et le vecteur propre \(W_c\) subsistent après compression. Le saut fini appartient déjà à l’espace invariant de jauge.\par
\Needspace{6\baselineskip}
\subsubsection{Unicité du processus reconstruit à partir de \(4\) positivités par réflexion}
La théorie euclidienne doit satisfaire la positivité par réflexion dans les quatre directions. Quatre mesures indépendantes ne sont pas construites. Une unique mesure de faces, obtenue à partir du même noyau réversible, factorise chaque forme d’Osterwalder–Schrader :\par
\[
\int
\overline{F(\Theta_{\nu}y)}F(y)\,d\Omega
=
\|\mathcal B_\nu F\|^2
\ge0,
\qquad \nu=1,\ldots,4.
\]
La positivité n’est pas une hypothèse ajoutée. C’est une norme au carré produite par la propriété de Markov et la réversibilité.\par
Les marginales des faces opposées retrouvent le même semi-groupe. Les quatre reconstructions OS possèdent donc une unique origine dynamique. Les connecteurs entre demi-réseaux démontreront ensuite qu’elles reconstruisent le même espace, le même vide et le même Hamiltonien.\par
\Needspace{6\baselineskip}
\subsubsection{Persistance du saut spectral dans la limite continue}
Les applications de raffinement entrelacent les formes de Dirichlet, les semi-groupes, le vide et le témoin \(W_c\). Le passage au continu ne repose pas seulement sur l’observation que le saut a la même valeur à chaque niveau. Cette égalité, isolément, n’empêcherait pas l’apparition de nouveaux états d’énergie arbitrairement faible à la limite.\par
La construction contrôle les produits croisés au moyen d’entrelaceurs de blocs. Les coordonnées grossières sont réparties sur leurs descendantes avec la normalisation exacte. Les identités de Gram démontrent que les expressions régularisées forment des familles de Cauchy. Ainsi, la forme, le vide et le vecteur propre atteignent la limite au sein du même système inductif.\par
Se obtiene\par
\[
D_{\infty,g}^{\rm YM}
\succeq
3\kappa_{\rm av}(I-P_{\Omega_\infty}),
\]
et le témoin limite satisfait\par
\[
D_{\infty,g}^{\rm YM}W_c
=
3\kappa_{\rm av}W_c.
\]
\Needspace{5\baselineskip}
Por tanto,\par
\[
\Delta_{\rm YM}^{\rm HMT}
=
3\kappa_{\rm av}>0.
\]
L’inégalité et le témoin convergent ensemble. Le saut n’est pas déduit d’une simple semi-continuité inférieure ni extrapolé à partir d’un maillage isolé.\par
\Needspace{6\baselineskip}
\subsubsection{Appartenance du témoin spectral au secteur de courbure}
Une objection subsisterait si \(W_c\) n’était qu’une excitation combinatoire d’un facteur interne, découplée de la courbure de Yang–Mills.\par
L’holonomie d’incidence élimine cette possibilité. La somme des six incidences détermine une holonomie locale d’ordre trois ; \(W_c\) est une fonction centrale de cette holonomie. Il appartient donc à l’algèbre locale invariante de jauge engendrée par les lacets.\par
Le lecteur \(F^2\) est polarisé pour produire un espace de courbures locales. Les entrelaceurs de raffinement transportent ces courbures en conservant le produit de \(\mathfrak g\). Les sommes cellulaires convergent vers une distribution covariante de jauge\par
\[
\mathbf F_{\mu\nu}
\in
\mathcal S'(\mathbb R^4;\mathfrak g)
\widehat\otimes
L^2(\widehat X_\infty,\widehat\mu_{\infty,g}^{\rm YM}).
\]
Le secteur cyclique engendré par \(\mathbf F\) est réducteur pour le Hamiltonien et contient le témoin \(W_c\). Le saut exact appartient donc au secteur physique de courbure. Ce n’est pas le spectre d’un bruit auxiliaire.\par
\Needspace{6\baselineskip}
\subsubsection{Localité euclidienne \(E(4)\) et reconstruction continue d’Osterwalder–Schrader}
Le générateur se décompose en circuits locaux de blocs à support uniformément borné. Une coloration nonadique sépare les colliers qui peuvent être mis à jour simultanément. Le produit des circuits préserve l’identité de Bianchi à chaque étape, et non seulement à la fin d’un balayage.\par
Les translations et rotations euclidiennes agissent par réétiquetage du réseau, des repères, des incidences et des holonomies. La continuité forte n’est pas prouvée uniquement pour un premier chaos abstrait. Elle est contrôlée d’abord sur les observables locales, puis sur les holonomies, les plaquettes et le trèfle ; la régularisation des orbites étend l’action à un domaine essentiel dense.\par
Le système cofinal produit un réseau local \(\mathfrak A_{\rm YM}(O)\), un vide \(\Omega_{\infty,g}\), une représentation forte de \(E(4)\) et des fonctionnelles de Schwinger. Les empilements de Markov et les transferts OS sont identifiés au même semi-groupe continu. La positivité par réflexion, la régularité, la covariance et la propriété de regroupement exponentiel permettent la reconstruction d’Osterwalder–Schrader.\par
Le prolongement produit des champs locaux de Wightman et une représentation de Poincaré dont le spectre est contenu dans le cône futur. Les quatre reconstructions euclidiennes coinduisent le même Hamiltonien de courbure.\par
Le saut contrôle en outre le regroupement :\par
\[
|S^{\rm YM}(A\,\tau_rB)|
\le
e^{-3\kappa_{\rm av}r}
\|A\Omega_\infty\|
\|B\Omega_\infty\|.
\]
Ainsi, la valeur spectrale positive a une conséquence physique continue et locale ; elle ne reste pas enfermée dans le calcul fini.\par
\Needspace{6\baselineskip}
\subsubsection{Reconstruction de l’action classique comme sortie du système euclidien}
La densité finie directrice est le lecteur exact de \(F^2\). Le trèfle ne s’y substitue pas. Il sert de carte lisse pour reconnaître sa limite.\par
Pour une connexion lisse, le développement de Magnus de l’holonomie d’une plaquette commence par\par
\[
\log\operatorname{Hol}_{p_{\mu\nu}}
=
a_m^2F_{\mu\nu}
+\text{terme impair}
+O(a_m^4).
\]
La moyenne sur les quatre orientations du trèfle annule le terme impair. La somme de Riemann donne\par
\[
S_{m,g}(A)
\longrightarrow
\frac1{4g^2}
\int_{\mathbb R^4}\|F_A(x)\|^2\,d^4x.
\]
La loi construite est donc reconnue comme Yang–Mills parce que son action finie, ses holonomies et sa courbure convergent vers la fonctionnelle classique correcte. On ne définit pas une « théorie de Yang–Mills » comme une théorie quelconque possédant le saut souhaité.\par
La même mesure satisfait Schwinger–Dyson et Ward. La relation entre petits lacets et courbure, l’identité de Bianchi et la covariance de jauge sont conservées à la limite. Ce n’est qu’après toute cette identification que l’équivalence spectrale transportant le saut est appliquée.\par
\Needspace{6\baselineskip}
\subsubsection{Interprétation dimensionnelle du paramètre de masse obtenu par l’entrelaceur terminal}
Le saut énergétique s’exprime dimensionnellement comme\par
\[
m_{\rm gap}^{\rm HMT}
=
\frac{\hbar_{\rm int}}{c_{\rm int}}
\,3\kappa_{\rm av}>0.
\]
Ce résultat n’attribue pas une masse de Proca à un gluon élémentaire. L’invariance de jauge demeure intacte. La masse correspond au coût minimal d’excitation du secteur collectif de courbure invariant de jauge au-dessus du vide.\par
Cette distinction ne diminue pas la portée du résultat : c’est exactement la formulation requise par le problème de Yang–Mills. Le saut appartient au spectre physique des états invariants de jauge, et non à un terme de masse élémentaire ajouté au lagrangien.\par
\Needspace{6\baselineskip}
\subsubsection{Existence d’un saut spectral positif et reconstruction massive du secteur de Yang–Mills}
Pour tout groupe de Lie compact, connexe et simple, la caractéristique de jauge du continu, suivie de son entrelaceur spécifique, construit une théorie quantique non triviale de Yang--Mills sur \(\mathbb R^4\).\par
La théorie possède :\par
\begin{itemize}
\item une famille projective de mesures locales dépendant du couplage ;
\item une action exacte reconnue à la limite comme \(\frac1{4g^2}\int\|F\|^2\) ;
\item un quotient de jauge physique ;
\item la positivité par réflexion dans quatre directions ;
\item un réseau local et une action forte de \(E(4)\) ;
\item des fonctions de Schwinger et un prolongement de Wightman ;
\item un vide simple ;
\item les équations de Schwinger–Dyson, de Ward et de Bianchi ;
\item un secteur cyclique de courbure ;
\item le saut exact
  \[
  \Delta_{\rm YM}^{\rm HMT}=3\kappa_{\rm av}>0.
  \]
\end{itemize}
La conclusion résiste aux objections prévisibles parce qu’elle n’identifie pas le générateur d’une fibre au produit ; ne laisse aucune coordonnée hors de l’indice ; ne confond pas la marginale des liens avec le quotient physique ; n’extrapole pas le saut à partir d’un maillage ; ne place pas le témoin hors du secteur de courbure ; ne présuppose pas la continuité euclidienne ; et n’utilise pas le saut pour construire l’action ou la mesure.\par
Navier–Stokes démontre que le continu peut transporter une non-linéarité sans perdre la mémoire nécessaire pour empêcher une singularité. Yang–Mills démontre que le même continu peut transporter une liberté de jauge sans perdre la courbure nécessaire pour séparer le vide de la première excitation physique.\par

\Needspace{6\baselineskip}
\section{Fermetures algébriques du continu et limites de la reconstruction extensionnelle}
Les cinq constructions \texttt{sp}, \texttt{sol}, \texttt{gau}, \texttt{coh} et \texttt{det} ne sont pas des branches ajoutées au continu pour résoudre cinq problèmes connus. Ce sont des caractéristiques simultanées du même état enrichi. Hodge et Birch–Swinnerton–Dyer sont les deux expressions finales de cette architecture : l’une transforme une histoire cohomologique compatible en classe algébrique ; l’autre démontre que les expressions analytique et arithmétique d’une même droite déterminantale coïncident.\par
\Needspace{6\baselineskip}
\subsection{4e résolution : Hodge et construction effective de la préimage algébrique}
Le problème de Hodge demande si toute classe rationnelle de type \((p,p)\),\par
\[
\alpha\in
H^{2p}(X,\mathbb Q)\cap H^{p,p}(X),
\]
pour une variété projective lisse complexe \(X\), provient d’une classe algébrique\par
\[
\beta_\alpha\in \operatorname{CH}^{p}(X)_{\mathbb Q}
\]
tal que\par
\[
\operatorname{cl}_{X,p}(\beta_\alpha)=\alpha.
\]
La difficulté ne consiste pas à écrire cette dernière égalité. Elle consiste à construire \(\beta_\alpha\) sans avoir préalablement introduit, sous un autre nom, l’existence de la préimage que l’on veut démontrer.\par
La caractéristique cohomologique \texttt{coh} fournit l’appareil antérieur à l’entrée de \(X\), \(p\) et \(\alpha\) : histoires survivantes, raffinements, relations, orientation, retenue, mémoire, frontière, parcours, multisection, terme terminal et lecteur. La donnée classique ne crée pas cet appareil ni ne sélectionne rétrospectivement un sous-domaine. Elle demande seulement ce qu’il exprime lorsqu’il reçoit les observations de \(\alpha\).\par
À chaque profondeur \(N\) apparaissent deux modules :\par
\[
\mathscr S_N(X,p),
\]
qui contient les états cohomologiques enrichis, et\par
\[
\mathscr H_N^p(X),
\]
qui contient les observations cohomologiques visibles à cette profondeur. Le lecteur fini\par
\[
e_N:\mathscr S_N(X,p)\longrightarrow\mathscr H_N^p(X)
\]
est coordonné aux applications de restriction :\par
\[
\rho_{N+1,N}:\mathscr S_{N+1}\longrightarrow\mathscr S_N,
\qquad
q_{N+1,N}:\mathscr H_{N+1}^p\longrightarrow\mathscr H_N^p,
\]
au moyen du carré\par
\[
e_N\rho_{N+1,N}
=
q_{N+1,N}e_{N+1}.
\]
Cela empêche qu’une observation admise à une profondeur soit incompatible avec son image au niveau précédent.\par
Pour la classe \(\alpha\), la fibre finie à peupler est\par
\[
\mathscr F_N(\alpha)
=
\{s\in\mathscr S_N(X,p):e_N(s)=q_N\alpha\}.
\]
La définir ne démontre pas qu’elle est non vide. Cela ne démontre pas non plus qu’un élément de \(\mathscr F_N(\alpha)\) peut se prolonger jusqu’à \(\mathscr F_{N+1}(\alpha)\). La résolution consiste précisément à construire ces deux propriétés.\par
\Needspace{6\baselineskip}
\subsubsection{Construction de la correction relative préalable à la classe de cohomologie}
Lors du raffinement de \(N\) à \(N+1\), deux noyaux relatifs apparaissent :\par
\[
\mathscr K_{N+1}^{\mathrm S}
=
\ker\rho_{N+1,N},
\qquad
\mathscr K_{N+1}^{\mathrm H}
=
\ker q_{N+1,N}.
\]
Le premier contient les corrections nouvelles de l’état qui disparaissent lorsqu’on revient au niveau précédent. Le second contient les observations nouvelles qui n’étaient pas non plus visibles au niveau précédent.\par
Les cartes nouvelles, les intersections de Mayer–Vietoris, les orientations, la retenue, la mémoire, la frontière, le parcours et les coordonnées terminales produisent une présentation finie. Sur ses générateurs sont construites deux applications :\par
\[
a_{N+1}:
\mathscr A_{N+1}^{\mathrm{rel}}
\longrightarrow
\mathscr K_{N+1}^{\mathrm S},
\]
\[
b_{N+1}:
\mathscr A_{N+1}^{\mathrm{rel}}
\longrightarrow
\mathscr K_{N+1}^{\mathrm H},
\]
satisfaisant\par
\[
e_{N+1}^{\mathrm{rel}}a_{N+1}=b_{N+1}.
\]
L’application \(a_{N+1}\) produit une correction de l’état ; \(b_{N+1}\) exprime exactement sa correction cohomologique.\par
Les observations relatives sont ordonnées selon la profondeur, le support, le feuillet, le parcours, la retenue, la mémoire et le terme terminal. Dans cet ordre, la matrice de \(b_{N+1}\) est unitriangulaire :\par
\[
b_{N+1}(c_\mu)
=
\bar c_\mu+
\sum_{\nu<\mu}
B_{\nu\mu}\bar c_\nu.
\]
La diagonale unité n’est pas choisie pour représenter une classe concrète. Elle provient de l’incidence principale de chaque générateur ; les termes inférieurs sont imposés par les relations, l’orientation et la retenue. Le tableau est construit sur toutes les observations relatives avant l’introduction de \(\alpha\).\par
L’élimination triangulaire produit explicitement un inverse à droite :\par
\[
\sigma_{N+1}^{\mathrm{rel}}(\bar c_{\mu_0})=c_{\mu_0},
\]
\[
\sigma_{N+1}^{\mathrm{rel}}(\bar c_\mu)
=
c_\mu-
\sum_{\nu<\mu}
B_{\nu\mu}
\sigma_{N+1}^{\mathrm{rel}}(\bar c_\nu),
\]
et, par récurrence,\par
\[
b_{N+1}\sigma_{N+1}^{\mathrm{rel}}=I.
\]
\Needspace{5\baselineskip}
Por tanto,\par
\[
s_{N+1}^{\mathrm{rel}}
=
a_{N+1}\sigma_{N+1}^{\mathrm{rel}}
\]
satisfait\par
\[
e_{N+1}^{\mathrm{rel}}
s_{N+1}^{\mathrm{rel}}
=I.
\]
C’est la protection contre la circularité. La section relative n’est pas supposée, ne se déduit pas du fait que la classe recherchée serait déjà algébrique et n’est pas construite spécifiquement pour \(\alpha\). Elle existe universellement pour toute observation relative.\par
Il n’est pas non plus nécessaire que la présentation soit injective. Le noyau peut contenir des expressions cellulaires redondantes. La preuve exige la surjectivité vers les observations relatives, et non l’unicité de l’expression. L’annulation du conoyau est une conséquence de l’inverse à droite ; ce n’est pas une condition directrice introduite avant la preuve.\par
\Needspace{6\baselineskip}
\subsubsection{L’opérateur universel d’extension}
La structure possède une dégénérescence unitaire\par
\[
\iota_{N+1,N}:\mathscr S_N\longrightarrow\mathscr S_{N+1},
\qquad
\rho_{N+1,N}\iota_{N+1,N}=I.
\]
Soit \(s\in\mathscr S_N\) et soit \(h\in\mathscr H_{N+1}^p\) une observation compatible :\par
\[
e_N(s)=q_{N+1,N}(h).
\]
Le défaut du prolongement trivial est\par
\[
\delta_{N+1}(s,h)
=
h-e_{N+1}\iota_{N+1,N}(s).
\]
La compatibilité garantit que\par
\[
\delta_{N+1}(s,h)
\in
\mathscr K_{N+1}^{\mathrm H}.
\]
On définit alors\par
\[
\operatorname{Ext}_{N+1,N}(s,h)
=
\iota_{N+1,N}(s)
+
s_{N+1}^{\mathrm{rel}}
\bigl(\delta_{N+1}(s,h)\bigr).
\]
Par construction :\par
\[
\rho_{N+1,N}
\operatorname{Ext}_{N+1,N}(s,h)
=s,
\]
\[
e_{N+1}
\operatorname{Ext}_{N+1,N}(s,h)
=h.
\]
La première identité conserve l’histoire antérieure. La seconde réalise exactement la nouvelle observation. La naturalité du tableau relatif garantit que l’orientation, la retenue, la mémoire, le parcours, la frontière et le terme terminal ne sont pas perdus lors du raffinement.\par
L’opérateur fonctionne pour tout couple compatible \((s,h)\). Il ne contient ni classe algébrique ni cycle choisi. C’est précisément pourquoi son application ultérieure à une classe de Hodge ne présuppose pas Hodge.\par
\Needspace{6\baselineskip}
\subsubsection{Construction niveau par niveau de l’histoire globale dans le système inverse}
À la profondeur initiale, il existe une section basée\par
\[
s_0^{\mathrm{base}}:
\mathscr H_0^p\longrightarrow\mathscr S_0.
\]
Ses images sont des états germes, non des cycles algébriques choisis pour représenter \(\alpha\). Ce n’est qu’ensuite que l’on fixe\par
\[
s_0=s_0^{\mathrm{base}}(q_0\alpha)
\]
et que l’on définit récursivement\par
\[
s_{N+1}
=
\operatorname{Ext}_{N+1,N}
\bigl(s_N,q_{N+1}\alpha\bigr).
\]
Les identités de l’opérateur universel prouvent par récurrence :\par
\[
e_Ns_N=q_N\alpha,
\qquad
\rho_{N+1,N}s_{N+1}=s_N.
\]
Chaque fibre \(\mathscr F_N(\alpha)\) est donc non vide et la famille\par
\[
\xi_\alpha=(s_N)_{N\ge0}
\]
est une histoire compatible :\par
\[
\xi_\alpha
\in
\varprojlim_N\mathscr F_N(\alpha).
\]
On n’invoque pas une limite inverse abstraite pour dissimuler l’absence de sections. La suite est présentée récursivement. L’annulation de \(\varprojlim^1\) n’est pas non plus utilisée en remplacement de la construction : la surjectivité et la condition de Mittag–Leffler apparaissent ensuite comme conséquences d’un opérateur d’extension déjà matérialisé.\par
\Needspace{6\baselineskip}
\subsubsection{Émergence de la classe algébrique comme limite de la construction relative}
Les observations finies et la coordonnée terminale séparent les expressions de la limite. Puisque \(\xi_\alpha\) reproduit \(q_N\alpha\) à toute profondeur,\par
\[
\operatorname{ev}_{\infty,X,p}^{\mathrm{Hdg}}
(\xi_\alpha)
=
\alpha.
\]
Ce n’est qu’alors que l’expression classique agit :\par
\[
X_\chi(X,p)
\xrightarrow{R_{\mathrm{Hdg}}}
K_0\!\left(
\operatorname{Perf}^{\ge p}(X)
\right)_{\mathbb Q}
\xrightarrow{\operatorname{ch}^{\mathrm{loc}}_p}
\operatorname{CH}^{p}(X)_{\mathbb Q}
\xrightarrow{\operatorname{cl}_{X,p}}
\operatorname{Hdg}^{p}(X)_{\mathbb Q}.
\]
Se define\par
\[
\beta_\alpha
=
\operatorname{ch}^{\mathrm{loc}}_p
\bigl(R_{\mathrm{Hdg}}(\xi_\alpha)\bigr).
\]
\Needspace{9\baselineskip}
L’identité d’action du lecteur,\par
\[
\operatorname{cl}_{X,p}
\operatorname{ch}^{\mathrm{loc}}_p
R_{\mathrm{Hdg}}
=
\operatorname{ev}_{\infty,X,p}^{\mathrm{Hdg}},
\]
da finalmente\par
\[
\operatorname{cl}_{X,p}(\beta_\alpha)
=
\alpha.
\]
Cette identité du lecteur ne prouve pas à elle seule l’existence de \(\xi_\alpha\). L’existence a été construite auparavant au moyen du tableau unitriangulaire, de l’inverse à droite, de l’opérateur universel et de la récurrence. Maintenir ces deux étapes séparées est précisément ce qui empêche la démonstration d’être circulaire.\par
Pour toute classe rationnelle de type \((p,p)\), la caractéristique cohomologique construit la section compatible \(\xi_\alpha\), le cycle \(\beta_\alpha\in\operatorname{CH}^p(X)_{\mathbb Q}\) et l’identité \(\operatorname{cl}_{X,p}(\beta_\alpha)=\alpha\), conformément à \cref{thm:pdf2-hodge-extension-universal,thm:pdf2-hodge-completitud-final}. L’exhaustivité cellulaire et FIN1 font partie de la construction démontrée de l’inverse à droite et de son prolongement coinductif ; ce ne sont pas des hypothèses externes de promotion.\par
\Needspace{6\baselineskip}
\subsection{5e résolution : BSD et coïncidence des expressions analytique et arithmétique}
La résolution de Birch–Swinnerton–Dyer ne commence pas par tenter de deviner l’ordre d’annulation de \(L(E,s)\), ni par réunir rétrospectivement rang, régulateur, groupe de Tate–Shafarevich et facteurs de Tamagawa.\par
La caractéristique déterminantale \texttt{det} conserve déjà histoire, produit fibré, unité commune, orientation, retenue, mémoire, survie, terme terminal, réseau, dualité, racine et expressions compatibles. Ce n’est qu’ensuite qu’intervient une courbe elliptique \(E/\mathbb Q\).\par
\Needspace{6\baselineskip}
\subsubsection{Complexe entier global et décomposition en fibres primaires}
La construction s’effectue sur un complexe libre entier d’histoires. Pour chaque nombre premier \(\ell\) et chaque puissance \(\ell^n\), les coefficients \(E[\ell^n]\), les transitions, Alexander–Whitney, le tiré en arrière de Manin, l’unité et les expressions sont obtenus par changement de base du même complexe entier.\par
La fibre ternaire ne gouverne pas les autres. On n’invente pas une fausse extension de \(\mathbb Z_3\) à \(\mathbb Z_\ell\). Cette précaution est indispensable pour que la conclusion englobe tout \(\Sha(E)\) et tous les facteurs de Tamagawa, et non uniquement sa composante \(3\)-primaire.\par
L’unité commune produit les expressions de Kato et de Poitou–Tate. Dans chaque bloc fini–singulier, on construit directement l’homotopie de remplissage\par
\[
h_*=-vf,
\]
dont le bord est la différence entre les deux expressions.\par
Les histoires sont ordonnées selon le nombre premier, le niveau, l’orientation, la retenue, la mémoire, la frontière, le terme terminal et la localisation. On obtient ainsi une décomposition basée exhaustive en :\par
\begin{itemize}
\item unité racine ;
\item blocs finis–singuliers ;
\item réseau local–global.
\end{itemize}
Aucune classe résiduelle ne demeure non typée.\par
L’homotopie \(H_{\mathrm{FS}}\) n’intervient pas comme prémisse. Dans chaque bloc normal, on définit son action et l’on vérifie\par
\[
\partial H_\lambda+H_\lambda\partial=I.
\]
La somme transportée par la décomposition basée satisfait\par
\[
\partial H_{\mathrm{FS}}
+
H_{\mathrm{FS}}\partial
=
p_{\mathrm{FS}}.
\]
Le membre de droite est le projecteur sur le secteur fini–singulier. Il n’efface pas le réseau qui conserve \(\Sha\) et Tamagawa. L’homotopie reproduit l’homotopie de remplissage préalablement construite ; elle n’est pas postulée pour faire disparaître une obstruction.\par
\Needspace{6\baselineskip}
\subsubsection{Finitude du groupe arithmétique préalable à l’évaluation du terme principal}
Le réseau résiduel possède une matrice entière \(D_E\). Avant de prendre son conoyau, on construit un inverse rationnel au moyen de la partie diagonale terminale et de la série géométrique finie de la partie strictement triangulaire. Les deux compositions sont vérifiées.\par
\Needspace{5\baselineskip}
Por tanto,\par
\[
D_E\otimes\mathbb Q
\]
est un isomorphisme et\par
\[
F_E=\operatorname{coker}D_E
\]
est fini. Cette finitude est obtenue avant toute utilisation du rang, de \(\Sha\), du régulateur ou du terme principal.\par
La forme de Smith conserve tous les diviseurs élémentaires et recompose exactement leurs composantes primaires. La finitude de \(\Sha\) n’est pas déduite de la formule finale : elle est dérivée ensuite de son inclusion dans un groupe fini déjà construit.\par
\Needspace{6\baselineskip}
\subsubsection{Bockstein, Kato–FERS et Poitou–Tate}
La différentielle déterminantale universelle \(S_E(T)\) engendre les pages de Bockstein. Sa forme de Smith \(T\)-adique détermine\par
\[
d
=
\operatorname{ord}_{T=0}
\det S_E(T)
=
\sum_q \dim V_q.
\]
Les pages ne sont pas remplacées par cet entier. Chaque quotient, chaque différentielle réduite, chaque accouplement et chaque jet sont conservés. Le coefficient régulier est\par
\[
R_{\mathrm{Boc}}^{\mathrm{der}}(S_E)
=
\left[
T^{-d}\det S_E(T)
\right]_{T=0}
\neq0.
\]
L’équivalence Kato–FERS est construite ligne par ligne. Alexander–Whitney coupe chaque histoire en tous ses points. La coupe terminale conserve le générateur de poids nul ; les autres coupes acquièrent des poids positifs déterminés par la profondeur et la retenue. Le tableau exhaustif prend la forme\par
\[
I+TB_i(T),
\]
avec un inverse convergent et un déterminant alterné égal à un dans \(T=0\). Ainsi, l’unité, les feuillets, la mémoire, l’orientation et le terme terminal sont préservés avant toute consultation d’une valeur de \(L(E,s)\).\par
Poitou–Tate relie ensuite les pages de Bockstein à Mordell–Weil. Les applications de Kummer, d’expression et de recollement sont écrites dans les deux sens et satisfont l’unité et la counité. Elles sont inverses avant la comparaison des dimensions.\par
Le dédoublement des jets conserve les \(q\) directions correspondant à une ligne de profondeur \(q\). On ne prétend pas obtenir \(q\) covecteurs à partir d’une seule coordonnée, ni compléter les dimensions par des identités artificielles. Le Gram conserve tous les termes croisés entre lignes, profondeurs et jets.\par
Ce n’est qu’ensuite que les dimensions sont prises :\par
\[
\operatorname{rank}E(\mathbb Q)
=
\sum_q\dim V_q
=
d.
\]
Le rang et l’ordre déterminantal coïncident parce que tous deux sont des expressions des pages construites, et non parce que l’un aurait été défini en recopiant l’autre.\par
\Needspace{6\baselineskip}
\subsubsection{Obtention du groupe de Tate–Shafarevich et des facteurs de Tamagawa par la suite exacte globale}
À partir du cône local réduit, on construit une injection\par
\[
\Sha(E)\hookrightarrow F_E.
\]
Les relèvements structurels des composantes locales produisent la projection et sa section. La vérification au niveau des cocycles donne :\par
\[
0
\longrightarrow
\Sha(E)
\longrightarrow
F_E
\longrightarrow
\bigoplus_v\Phi_v(\mathbb F_v)
\longrightarrow
0.
\]
Puisque \(F_E\) était déjà fini,\par
\[
|\Sha(E)|<\infty
\]
y\par
\[
|F_E|
=
|\Sha(E)|
\prod_v c_v.
\]
L’exactitude n’est pas déduite d’une comparaison des cardinaux. Les applications sont d’abord construites, leurs noyaux et leurs images identifiés ; ce n’est qu’ensuite que leurs ordres sont pris.\par
\Needspace{6\baselineskip}
\subsubsection{Construction indépendante du membre analytique de BSD}
La modularité de \(E\) fournit la forme nouvelle associée. Sa transformée de Mellin est divisée en deux intervalles, et l’équation fonctionnelle transforme la seconde moitié en une intégrale convergente avec toutes ses dérivées. Cela produit un germe holomorphe de \(L(E,s)\) autour de \(s=1\) sans supposer son ordre d’annulation.\par
Dans le demi-plan de convergence absolue, les facteurs d’Euler sont réalisés par des opérateurs normaux. Les troncatures convergent en norme trace et leur déterminant de Fredholm retrouve le produit d’Euler. La réduction de Schur sur les coordonnées terminales et un atlas de cartes déterminantales transportent cette famille jusqu’à la différentielle \(S_E\).\par
Il en résulte une factorisation construite :\par
\[
L(E,s)
=
u_E(s)
\det S_E(T_E(s)).
\]
\Needspace{7\baselineskip}
La carte satisfait\par
\[
T_E(1)=0,
\qquad
T'_E(1)=1,
\]
et l’unité de transition satisfait\par
\[
u_E(1)=1.
\]
Cette normalisation est démontrée au moyen des transitions de l’atlas, du complément de Fredholm et des changements de base. Elle n’est pas ajustée au moyen du rang, du régulateur, de \(\Sha\) ni du coefficient final.\par
En consecuencia,\par
\[
L(E,s)
=
(s-1)^dU_E^{\mathrm{lead}}(s),
\]
con\par
\[
U_E^{\mathrm{lead}}(1)
=
R_{\mathrm{Boc}}^{\mathrm{der}}(S_E)
\neq0.
\]
\Needspace{5\baselineskip}
Por tanto,\par
\[
\operatorname{ord}_{s=1}L(E,s)=d.
\]
En combinant cela avec Poitou–Tate :\par
\[
\operatorname{ord}_{s=1}L(E,s)
=
\operatorname{rank}E(\mathbb Q).
\]
\Needspace{6\baselineskip}
\subsubsection{Dérivation interne du régulateur à partir de l’accouplement des hauteurs}
La base entière de Mordell–Weil provient du système complet de jets. Sur elle sont construites la hauteur de Néron–Tate et sa décomposition locale–globale.\par
Le Gram contient tous les termes croisés :\par
\[
G_E^{\mathrm{jet}}
=
\bigl(
\langle P_{q,a},P_{q',a'}\rangle_{\mathrm{NT}}
\bigr).
\]
On ne suppose pas de diagonalité entre profondeurs et l’on ne remplace pas les directions supplémentaires par des blocs identité. Son déterminant définit\par
\[
\operatorname{Reg}(E)
=
\det G_E^{\mathrm{jet}}
>0.
\]
L’extension des scalaires à \(\mathbb R\) est effectuée avant la multiplication par le régulateur. Le comparateur entre les droites de Bockstein et de Néron est donc bien typé ; on ne tente pas d’introduire un nombre réel dans une droite encore définie uniquement sur \(\mathbb Q\).\par
\Needspace{6\baselineskip}
\subsubsection{Diagramme de \(4\) morphismes et évaluation du terme principal}
La comparaison finale utilise quatre flèches de droites déterminantales :\par
\begin{enumerate}
\item Kato–FERS transporte la section analytique.
\item Poitou–Tate la conduit jusqu’à Mordell–Weil et à sa hauteur.
\item La suite exacte sépare \(\Sha\) et Tamagawa.
\item L’évaluation intègre la torsion, la période de Néron et les composantes locales.
\end{enumerate}
Chaque flèche possède un domaine, un codomaine et une action sur une base. Avant l’introduction de l’élément analytique, deux trivialisations indépendantes sont fixées : l’analytique et celle de Néron.\par
La commutativité du carré final provient de la conservation des bases structurelles. Elle n’est pas imposée en égalisant préalablement les nombres qui doivent apparaître.\par
L’application de la composition à l’élément analytique donne :\par
\[
\operatorname{ord}_{s=1}L(E,s)
=
\operatorname{rank}E(\mathbb Q)
=
r,
\]
y\par
\[
\frac{L^{(r)}(E,1)}
{r!\,\Omega_E}
=
\frac{
|\Sha(E)|
\operatorname{Reg}(E)
\prod_v c_v
}{
|E(\mathbb Q)_{\mathrm{tors}}|^2
}.
\]
C’est la formule complète de Birch--Swinnerton--Dyer pour toute courbe elliptique \(E/\mathbb Q\), comme l’établit \cref{thm:pdf2-bsd-publicacion-final}. La famille compatible \(\nu_{E,\ell,n}\) est construite au sein du complexe entier global et fournit le comparateur qui relie ses fibres primaires ; elle fait partie de la démonstration universelle et ne constitue pas une condition externe de promotion.\par
Sa protection contre la circularité est strictement chronologique :\par
\begin{itemize}
\item le complexe entier est fixé avant la séparation des nombres premiers ;
\item l’homotopie de remplissage est construite avant \(H_{\mathrm{FS}}\) ;
\item Smith est obtenu avant le rang ;
\item Kato–FERS est vérifié avant le déterminant principal ;
\item Poitou–Tate est construit avant la comparaison des dimensions ;
\item l’exactitude \(\Sha\)–Tamagawa est prouvée avant de prendre les ordres ;
\item le prolongement Mellin–Fredholm est construit avant l’ordre analytique ;
\item la hauteur est construite avant le régulateur ;
\item les quatre flèches sont fixées avant l’introduction de l’élément analytique ;
\item les trivialisations sont fixées avant l’évaluation des deux membres.
\end{itemize}
Aucun élément n’est normalisé au moyen de la valeur qu’il doit ensuite démontrer.\par
\Needspace{6\baselineskip}
\subsection{Factorisation des entrelaceurs terminaux RH–NS–YM–Hodge–BSD à travers le continu généalogique commun}
Riemann, Navier–Stokes, Yang–Mills, Hodge et BSD ne sont pas cinq traductions arbitraires d’un même vocabulaire. Chaque problème apparaît lorsqu’une expression classique oublie une partie différente de l’état enrichi :\par
\begin{itemize}
\item Riemann oublie la colligation commune qui conserve simultanément les canaux premier, gamma et polaire.
\item Navier–Stokes oublie la colonne causale complète comprenant terme terminal, pertes, retenue et mémoire.
\item Yang–Mills oublie la coordination entre incidence, courbure, loi, dynamique et secteur physique.
\item Hodge conserve une classe cohomologique, mais pas l’histoire compatible qui produit sa préimage algébrique.
\item BSD sépare l’expression analytique de l’arithmétique et perd la droite déterminantale commune qui les transporte.
\end{itemize}
L’architecture HMT n’impose pas cinq solutions indépendantes. Elle conserve d’abord ce que les formulations classiques expriment séparément :\par
\[
\begin{aligned}
\text{état enrichi}
&\longrightarrow \text{raffinements compatibles}
\longrightarrow \text{survie et terminal}\\
&\longrightarrow \text{continu discret}
\longrightarrow \text{publications classiques}.
\end{aligned}
\]
La raison pour laquelle les solutions fonctionnent n’est pas une similitude verbale. Dans les cinq cas, le théorème décisif provient d’une identité de conservation matérielle :\par
\begin{itemize}
\item conservation de l’énergie ou du défaut ;
\item conservation de l’histoire sous raffinement ;
\item conservation de l’action et du produit surfacique ;
\item conservation de l’observation durant l’extension ;
\item conservation de l’unité et de la base dans les droites déterminantales.
\end{itemize}
C’est pourquoi les problèmes classiques doivent apparaître à la fin de l’organisation éditoriale. L’objet commun est construit en premier ; HMT–MD développe ensuite ses conséquences ; la Mécanique dimensionnelle réalise alors ses transducteurs ; ce n’est qu’en dernier lieu que les cinq conclusions sont lues dans leurs langages classiques.\par
\Needspace{6\baselineskip}
\subsection{Distinction typée entre les cinq fermetures terminales et le corridor logique ZFC--cycle}
La relation avec Zermelo--Fraenkel, ZFC et la construction de von Neumann doit être formulée avec précision et dans l’ordre causal du corpus. Les cinq résolutions terminales et le corridor logique ZFC--cycle sont des démonstrations distinctes, fondées sur la même structure discrète du continu : la portée locale de l’une ne peut servir de négation transversale de l’autre. Les cinq fermetures établissent leurs résultats classiques respectifs ; le corridor logique spécifique conserve en outre la chaîne\par
\[
\begin{aligned}
 \mathcal C_{\mathrm{cont}}^{\mathrm{disc}}
 &\longrightarrow\mathcal L_0\xrightarrow{\Phi}\mathcal L_1
 \longrightarrow\gamma_F\longrightarrow w(\gamma_F)\\
 &\longrightarrow\mathsf{NoEscape}_{24}
 \longrightarrow\mathsf{AntiDiag}_{369}\\
 &\longrightarrow\mathsf{G\ddot{o}delIII}_{\mathrm{HMT}}
 \longrightarrow\mathsf{CycleCriterion}_{\mathrm{ZFC}},
\end{aligned}
\]
et exprime le critère fort retrouvé\par
\[
 \boxed{
 \neg\operatorname{Consis}(\mathrm{ZFC})
 \Longleftrightarrow
 \exists C\;\bigl(C\text{ est un cycle hors bande et }|C|>24\bigr).}
\]
Cette équivalence possède son propre domaine, ses encodeurs et sa preuve ; elle ne se déduit pas des cinq théorèmes terminaux et n’est pas limitée par une observation locale concernant l’un d’eux. De même, la résolution cardinale déjà construite conserve simultanément\par
\[
 \mathsf{CH}_{\mathrm{HMT}},
 \qquad
 M_h=L[h]\models 2^{\aleph_0}=\aleph_1,
\]
sans transformer la droite réelle ou l’énoncé extensionnel en entrées du générateur.\par

Une représentation extensionnelle formulée dans le langage de ZFC peut coder les supports ensemblistes des objets utilisés :\par
\begin{itemize}
\item un état fini est un ensemble ;
\item une histoire est une suite ;
\item un arbre est un ensemble de mots ;
\item une transition est une fonction ou une relation ;
\item un système inverse est une famille ensembliste ;
\item sa limite est un sous-ensemble d’un produit.
\end{itemize}
Ce codage ultérieur ne contient pas à lui seul l’identité généalogique de l’objet et ne gouverne pas les théorèmes construits avant l’expression extensionnelle. L’insuffisance apparaît lorsque le réduit extensionnel est traité comme s’il épuisait le parcours, l’orientation, la retenue, la mémoire, la frontière et le terme terminal.\par
L’axiome d’extensionnalité affirme :\par
\[
\forall A\,\forall B
\left[
\forall x\,
(x\in A\leftrightarrow x\in B)
\Rightarrow A=B
\right].
\]
L’assertion est correcte pour les ensembles. Deux ensembles ayant les mêmes éléments sont le même ensemble.\par
Mais il n’en découle pas qu’une expression extensionnelle permette de reconstruire l’histoire oubliée lors de sa production.\par
Sea\par
\[
\operatorname{Pub}:
\mathscr X_{\mathrm{gen}}
\longrightarrow
\mathscr X_{\mathrm{ext}}
\]
\Needspace{6\baselineskip}
un lecteur qui envoie un état généalogique sur sa valeur extensionnelle. Il peut arriver que\par
\[
\operatorname{Pub}(x)
=
\operatorname{Pub}(y)
\]
tandis que\par
\[
x\neq y
\]
en tant qu’histoires, parce qu’elles possèdent des parcours, orientations, retenues, mémoires, frontières ou termes terminaux différents.\par
Le réduit identifie les deux images dans \(\mathscr X_{\mathrm{ext}}\), mais cette égalité dans le codomaine n’implique pas \(x=y\) dans le domaine enrichi. L’erreur de type consiste à transformer l’égalité des expressions en égalité généalogique.\par
On peut définir la relation\par
\[
x\sim_{\mathrm{ext}}y
\quad\Longleftrightarrow\quad
\operatorname{Pub}(x)
=
\operatorname{Pub}(y).
\]
Alors, dans le domaine correspondant,\par
\[
\mathscr X_{\mathrm{ext}}
\simeq
\mathscr X_{\mathrm{gen}}/
{\sim_{\mathrm{ext}}}.
\]
La fibra\par
\[
\operatorname{Pub}^{-1}(a)
\]
contient les généalogies qui produisent la même image \(a\). L’extensionnalité décrit correctement le quotient ; elle ne fournit pas une section canonique reconstruisant une histoire concrète de chaque fibre.\par
Même l’axiome du choix ne résout pas cette carence. Le choix peut sélectionner un représentant de chaque fibre sous ses hypothèses, mais :\par
\begin{itemize}
\item il ne le rend pas canonique ;
\item il ne conserve pas nécessairement les opérations ;
\item il ne garantit pas la naturalité sous raffinement ;
\item il ne reconstruit ni retenue, ni mémoire, ni terme terminal ;
\item il ne démontre pas que la sélection correspond à la généalogie physique ou mathématique correcte.
\end{itemize}
Choisir un représentant n’équivaut pas à retrouver sa provenance.\par
\Needspace{6\baselineskip}
\subsubsection{L’exemple des ordinaux de von Neumann}
La construction de von Neumann\par
\[
0=\varnothing,
\qquad
n+1=n\cup\{n\}
\]
est exacte. L’ordinal\par
\[
n=\{0,1,\ldots,n-1\}
\]
code parfaitement son ordre extensionnel.\par
\Needspace{6\baselineskip}
Mais le même nombre peut apparaître comme sortie d’histoires différentes :\par
\[
6=3+3=2\cdot3=7-1.
\]
En tant qu’ordinal de von Neumann, le résultat est le même ensemble. C’est correct : l’arithmétique exige que les quatre expressions produisent le même \(6\).\par
Ce que l’ordinal seul ne conserve pas, c’est l’opération, le parcours, la retenue ou l’orientation qui a produit cette occurrence concrète. Si cette provenance importe pour une application ultérieure, elle doit être maintenue comme partie de l’objet enrichi :\par
\[
\widehat n_x
=
\{(k,x_k):k<n\}.
\]
La projection\par
\[
(k,x_k)\longmapsto k
\]
retrouve exactement l’ordinal de von Neumann. HMT n’élimine pas l’ordinal ; elle le reconnaît comme expression d’une structure susceptible de contenir davantage d’information.\par
La hiérarchie cumulative \(V_\alpha\) conserve l’appartenance et le rang. Elle n’enregistre pas automatiquement une causalité interne de construction différente du codage choisi. Réintroduire des étiquettes, des arbres ou des catégories permet de transporter des données généalogiques supplémentaires, mais cet enrichissement confirme précisément que la généalogie n’était pas contenue dans la valeur extensionnelle après avoir été oubliée.\par
\Needspace{6\baselineskip}
\subsubsection{La droite réelle comme expression ultérieure}
La même distinction concerne le continu. Un nombre réel peut être construit comme coupure de Dedekind, classe de suites de Cauchy ou autre représentant extensionnel équivalent. Cette équivalence est mathématiquement correcte.\par
HMT ajoute une question préalable : quel système compatible d’états, de raffinements et de termes terminaux exprime ce réel ?\par
La recta\par
\[
\mathbb R
\]
apparaît ensuite par une évaluation archimédienne. Elle conserve la valeur exprimée, mais pas nécessairement l’histoire complète qui l’a produite. Deux états peuvent partager la même évaluation réelle sans partager la même généalogie.\par
La structure discrète du continu ne doit donc pas introduire \(\mathbb R\) comme entrée. L’histoire compatible est construite en premier ; son image archimédienne est exprimée ensuite.\par
L’entrelaceur avec la cardinalité classique est construit explicitement. Le codage par blocs fournit l’homéomorphisme \(\beta:\Omega_\Gamma^{\rm cal}\simeq(\mathbb F_3^6)^{\mathbb N}\) ; ses approximants rationnels construisent \(\mathbb R_{\rm HMT}\), et le quotient par égalité de valeur exprime \(\mathbb R\). Le code généalogique \(h\) induit le bon ordre et détermine le modèle canonique \(M_h=L[h]\), dans lequel la construction prouve l’énoncé classique \(2^{\aleph_0}=\aleph_1\). Les résultats d’indépendance de Gödel et de Cohen classifient le langage extensionnel de ZFC après oubli de \(h\) ; ils ne réfutent pas la décision produite par la structure enrichie qui conserve le sélecteur.\par
\Needspace{6\baselineskip}
\subsubsection{Distinction entre identité généalogique, équivalence observationnelle et égalité extensionnelle}
La carence est typologique : un codage ensembliste ajouté peut transporter une généalogie déjà construite, tandis que le réduit extensionnel nu est dépourvu de l’application qui a produit l’expression. Capacité de codage et suffisance fondationnelle sont des assertions différentes.\par
Trois notions doivent être distinguées :\par
\[
\text{identité généalogique},
\qquad
\text{équivalence observationnelle},
\qquad
\text{égalité extensionnelle}.
\]
La première exige l’égalité de l’histoire, du parcours, de la mémoire, de la retenue et du terme terminal. La deuxième exige l’égalité relativement à un lecteur déterminé. La troisième exige l’égalité de l’objet exprimé dans le codomaine ensembliste.\par
Les confondre engendre de faux problèmes de reconstruction. Une image peut être exacte et, en même temps, insuffisante pour retrouver l’objet qui l’a produite.\par
C’est l’apport fondationnel que les cinq résolutions rendent visible. Chaque problème classique travaillait avec une expression correcte, mais appauvrie :\par
\begin{itemize}
\item une forme explicite sans toute sa généalogie commune ;
\item un champ sans mémoire causale complète ;
\item une théorie de jauge sans coordination intégrale de son processus ;
\item une classe cohomologique sans son histoire d’extension ;
\item une fonction \(L\) séparée de sa droite arithmétique.
\end{itemize}
La structure HMT maintient ces corrélations avant d’en exprimer les images. Elle n’a donc pas à reconstruire a posteriori, dans chaque problème, une information qui n’a jamais été détruite.\par
La conclusion fondationnelle est la suivante :\par
\begin{quote}
La fermeture extensionnelle est une expression ultérieure de la fermeture généalogique et ne peut s’y substituer. HMT construit la structure enrichie avant la droite réelle, démontre \(\mathsf{CH}_{\mathrm{HMT}}\), réalise la fermeture ordinale dans \(L[h]\) et établit séparément le corridor Gödel--HMT III et le critère ZFC--cycle. Un codage extensionnel ultérieur ne conserve que ce que déclare le morphisme d’expression ; il ne révoque aucun de ces résultats.\par
\end{quote}
\Needspace{6\baselineskip}
\subsection{Factorisation de RH–NS–YM–Hodge–BSD comme expressions terminales du continu généalogique commun}
Les cinq problèmes apparaissent ainsi comme cinq lectures terminales d’un unique bord. Ils ne sont pas résolus parce qu’une analogie les réunirait, mais parce que le continu enrichi conserve dès l’origine les corrélations que chaque formulation classique avait séparées.\par
La dépendance causale définitive se factorise comme suit :\par
\[
\begin{aligned}
\mathrm{APP}
&\longrightarrow \mathrm{TRIT}
\longrightarrow \mathrm{TPK}
\longrightarrow \widehat x\\
&\longrightarrow \text{cinq caractéristiques consubstantielles}
\longrightarrow \text{terminal et limite}
\longrightarrow \mathfrak C_{\mathrm{cont}}^{\mathrm{disc}}
\end{aligned}
\]
et seulement ensuite :\par
\[
\mathfrak C_{\mathrm{cont}}^{\mathrm{disc}}
\longrightarrow
\begin{cases}
\text{positivité de Weil et RH},\\
\text{régularité globale lisse de Navier--Stokes pour toute donnée périodique lisse},\\
\text{Yang--Mills et saut de masse pour tout groupe compact, connexe et simple},\\
\text{algébricité de toute classe rationnelle de Hodge},\\
\text{égalité complète de BSD pour toute courbe elliptique sur }\mathbb Q.
\end{cases}
\]
Les cinq lignes finales ne sont pas cinq branches internes du continu. Ce sont cinq expressions classiques ultérieures de cinq capacités indissociables qui coexistaient déjà dans le même état.\par
Telle est la fermeture que la formulation précédente ne pouvait pas encore formuler : il ne s’agit pas de cinq problèmes déconnectés qui admettraient par hasard un même vocabulaire, mais de cinq pertes d’information distinctes corrigées par une seule architecture qui a refusé de la détruire.\par
